\paraid{p-96eb53c66867}
In Part~\ref{part:topology},
we proved that the Gromov--Hausdorff space
is an absolute retract.
We now prove the discrete approximation property needed to identify
this space with Hilbert space.
Let
$\GHSpace$
denote the space of isometry classes of nonempty compact metric spaces
with the ordinary unpointed Gromov--Hausdorff distance
$\GHDistance$.
Antonyan asks whether
$\GHSpace$
is homeomorphic to the separable
infinite-dimensional Hilbert space
\cite[p.~2]{Antonyan2020Euclidean}.
For each positive integer
$\FiniteSize $,
Antonyan identifies the subspace represented by compact subsets of
$\RealNumbers^\FiniteSize $
with the Hilbert cube minus a point
\cite[Corollary~5.3]{Antonyan2021Euclidean}.

\paraid{p-707b7043674a}
Previous work has explored the topology of $\GHSpace$ through explicit
families and embeddings.
The author constructed families of branching Gromov--Hausdorff geodesics
continuously parametrized by the Hilbert cube
\cite[Theorem~1.3]{Ishiki2022Branching}.
The author also constructed topological embeddings of arbitrary compact metrizable
spaces into the subsets of
$\GHSpace$
consisting of continua
\cite[Theorem~1.1]{Ishiki2022Continua},
spaces with prescribed dimensions
\cite[Theorem~1.3]{Ishiki2023FractalDimensions},
and compact metric trees
\cite[Theorem~1.1]{Ishiki2023MetricTrees}.
In each construction,
finitely many distinct values of the embedding may be prescribed.
Byakuno
\cite[arXiv v1, Corollary~1.4]{Byakuno2026Embedding}
proves that countable products of closed intervals with positive summable lengths,
equipped with the supremum metric,
admit isometric embeddings into
$\GHSpace$.

\paraid{p-1084d3b137e6}
We also recall the dimension of the subspace of finite metric spaces.
Nakajima,
Yamauchi,
and Zava
\cite[arXiv v1, Theorem~A]{NakajimaYamauchiZava2025}
prove that the subspace of spaces with at most
$\FiniteSize$
points has topological dimension
$\FiniteSize(\FiniteSize-1)/2$.
Nakajima and Shioya
\cite[arXiv v1, Main Theorem~1.2]{NakajimaShioyaCompactification}
construct a compact metrizable space containing a dense topological copy of
$\GHSpace$.

\paraid{p-c46ed027b219}
A related Hilbert space classification concerns the Urysohn universal
metric space.
Uspenskij
\cite[arXiv v1, Theorem~2.1]{Uspenskij2004Urysohn}
proves that this space is homeomorphic to
$\SeparableHilbertSpace$.

\paraid{p-6d7aa84e3e0d}
Our approximation theorem concerns a countable family of prescribed maps
from compact metrizable domains.
We approximate all these maps so that every point of
$\GHSpace$
has a neighborhood meeting at most one of their images.
The construction uses products with finite equilateral spaces,
independently of Parts~\ref{part:measures}--\ref{part:topology}.
We combine this discrete approximation with the absolute retract property
to prove the following theorem.

\begin{theorem}\label{thm:part-iv-main}
\paraid{p-13f85da9e5b7}
The space
$\GHSpace$
is homeomorphic to the real Hilbert space
$\SeparableHilbertSpace$.
\end{theorem}

\paraid{p-add73a912728}
The classification concerns only topology and does not preserve the geometric
information carried by the Gromov--Hausdorff distance.
Ivanov and Tuzhilin
\cite[Main Theorem]{IvanovTuzhilin2019Isometries}
prove that every surjective isometry from
$\GHSpace$ to itself
is the identity.
Zava
\cite[arXiv v3, Theorem~B]{Zava2025Coarse}
proves that the subspace of finite metric spaces cannot be coarsely
embedded into any Hilbert space.
Consequently,
a homeomorphism in \Cref{thm:part-iv-main} cannot be bi-Lipschitz.

\paraid{p-60f77024d745}
For a sequence of compact metrizable domains
$\{\ApproximationDomain_\FamilyIndex\}_{\FamilyIndex\in\NonnegativeIntegers}$,
continuous maps
$\InputMap_\FamilyIndex\colon\ApproximationDomain_\FamilyIndex\to\GHSpace$,
and an open cover
$\OpenCover$
of
$\GHSpace$,
we construct continuous maps
$\ApproximatingMap_\FamilyIndex\colon\ApproximationDomain_\FamilyIndex\to\GHSpace$
such that
\[
 \forall\FamilyIndex\in\NonnegativeIntegers\
 \forall\DomainPoint\in\ApproximationDomain_\FamilyIndex\
 \exists U\in\OpenCover,\qquad
 \{\InputMap_\FamilyIndex(\DomainPoint),
   \ApproximatingMap_\FamilyIndex(\DomainPoint)\}\subset U,
\]
and
\[
 \{\ApproximatingMap_\FamilyIndex(\ApproximationDomain_\FamilyIndex)
   \}_{\FamilyIndex\in\NonnegativeIntegers}
 \text{ is discrete in }\GHSpace.
\]
The first condition means that the maps are
$\OpenCover$-close.
The second means that every point of
$\GHSpace$
has a neighborhood meeting at most one image.
This is \Cref{thm:discrete-approximation}.
Taking
$\ApproximationDomain_\FamilyIndex=\HilbertCube
 =\UnitInterval^{\NonnegativeIntegers}$
for every
$\FamilyIndex$
yields the approximation
condition in Toru\'nczyk's characterization,
stated in \Cref{thm:torunczyk}
and taken from
\cite[p.~248, assertion~(i)]{Torunczyk1981},
with the correction in
\cite[Section~C]{Torunczyk1985}.
Together with completeness,
separability,
and the absolute retract property,
this condition proves \Cref{thm:part-iv-main}.

\paraid{p-01c02735fab9}
The main theorem also implies that every point has a basis of open neighborhoods
that strongly deformation retract onto that point.
We do not know whether these neighborhoods can be chosen to be balls for
$\GHDistance$.

\medskip
\noindent\textbf{Dependence on the other parts.}
\paraid{p-518fca37991d}
\Cref{thm:discrete-approximation} uses the common preliminaries
and the metric constructions recalled in Section~\ref{sec:prelim-hilbert}.
The absolute retract result \Cref{thm:part-iii-main} enters only when we
apply Toru\'nczyk's characterization in
Subsection~\ref{subsec:approximation-hilbert-recognition}.
Thus the proof of the classification combines the approximation in this
part with the result of Parts~\ref{part:measures}--\ref{part:topology}.

\medskip
\noindent\textbf{Organization.}
\paraid{p-c7e99d59b5f0}
In Section~\ref{sec:prelim-hilbert},
we define discrete families and state the Hilbert space characterization.
In Subsection~\ref{subsec:approximation-products},
we recall the Gromov--Hausdorff estimate for metric products.
In Subsection~\ref{subsec:approximation-discrete},
we construct approximations of maps from countably many
compact domains whose images form a discrete family,
proving \Cref{thm:discrete-approximation}.
In Subsection~\ref{subsec:approximation-hilbert-recognition},
we apply Toru\'nczyk's characterization to prove
\Cref{thm:part-iv-main} and obtain strong local contractions.
In \Cref{sec:questions},
we conclude with questions about related spaces
and the geometry of Gromov--Hausdorff balls.

\medskip
\noindent\textbf{Conventions and notation.}
\paraid{p-bb6d5f435441}
All compact metric spaces are nonempty,
and Banach and Hilbert spaces are real.
We identify compact metric spaces with their isometry classes in
$\GHSpace$.
An isometry is a surjective isometric embedding.
For a metric space
$\MetricSpace{\BaseCarrier}{\MetricSymbol}$,
we write
$\CompactHyperspace(\BaseCarrier)$
for its nonempty compact subsets,
equipped with the Hausdorff metric
$\HausdorffDistance{\MetricSymbol}$.
For a compact space with an isometric embedding into an ambient metric space,
we identify the compact space with its image and use the restricted ambient metric.
The notation
$\IdentityMap$
denotes the identity map on its specified domain.
Sequences are indexed by
$\NonnegativeIntegers=\{0,1,2,\ldots\}$.
